\documentclass[../../../include/open-logic-section]{subfiles}

\begin{document}

\olfileid{sth}{z}{infinity-again}	
\olsection{اللانهاية}

لدينا بالفعل ما يكفي من المسلّمات لضمان وجود عدد لا نهائي من المجموعات
(إذا وجدت أي مجموعة). افترض مثلًا أن مجموعة ما موجودة، ومن ثم توجد
$\emptyset$ (بموجب~\olref[sth][z][sep]{prop:emptyexists}). ولكل
مجموعة~$x$، توجد المجموعة~$x \cup \{x\}$ بموجب
\olref[sth][z][pairs]{prop:pairsconsequences}. ولذلك، بتطبيق هذه العملية
بضع مرات، نحصل على المجموعات الآتية:
\begin{enumerate}
	\item[0.] $\emptyset$ %& $0$ & stage $0$\\
	\item[1.] $ \{\emptyset\}$ %& $1$ & stage $1$\\
	\item[2.] $\{\emptyset, \{\emptyset\}\}$ %& $2$ & stage $2$\\
	\item[3.] $\{\emptyset, \{\emptyset\}, \{\emptyset, \{\emptyset\}\}\}$ %&$3$ & stage $3$\\
	\item[4.] $\{\emptyset, \{\emptyset\}, \{\emptyset, \{\emptyset\}\}, \{\emptyset, \{\emptyset\}, \{\emptyset, \{\emptyset\}\}\}\}$% & $4$ & stage $4$ 
\end{enumerate}
ويمكننا التحقق من أن كلًا من هذه المجموعات متميز.

بدأنا الترقيم من~$0$ لأسباب عدة. وأحدها أنه ليس من الصعب جدًا التحقق
من أن المجموعة التي سميناها «$n$» لها~$n$ من العناصر بالضبط، وأنها
(بحسب الفهم الحدسي) تتكون في المرحلة رقم~$n$.

لكن هذا يمنحنا \emph{عددًا لا نهائيًا من} المجموعات، ولا يضمن وجود
\emph{مجموعة لا نهائية}، أي مجموعة لها عدد لا نهائي من العناصر.
وهذا مهم حقًا: فإذا لم نجد مجموعة لا نهائية (على معنى ددكند)، لم
نستطع بناء جبر ددكندي. لكننا نريد جبرًا ددكنديًا لكي نعامله بوصفه
مجموعة الأعداد الطبيعية. (قارن
\olref[sfr][infinite][dedekindsproof]{sec}.)

والأهم أن المسلّمات التي وضعناها إلى الآن \emph{لا} تضمن وجود أي
مجموعة لا نهائية. ولذلك علينا أن نضع مسلّمة جديدة:

\begin{axiom}[اللانهاية]
توجد مجموعة~$I$ بحيث~$\emptyset \in I$، ويكون
$x \cup \{x\} \in I$ كلما كان~$x \in I$.
\begin{align*}
	\exists I( & (\exists o \in I)\forall x\ x \notin o \land {}\\
	& (\forall x \in I)(\exists s \in I)\forall z(z \in s \liff (z \in x \lor z = x)))
\end{align*}
\end{axiom}

من السهل رؤية أن المجموعة~$I$ التي تمنحنا إياها مسلّمة اللانهاية
لا نهائية على معنى ددكند. وعنصرها المميز هو~$\emptyset$، والحقن على
$I$ معطى بـ$s(x) = x\cup\{x\}$. وقد بيّنت
\olref[sfr][infinite][dedekind]{thm:DedekindInfiniteAlgebra} كيفية
استخراج جبر ددكندي من مجموعة لا نهائية على معنى ددكند؛ وسنعامل هذا
الجبر بوصفه مجموعة الأعداد الطبيعية. وبتعبير أدق:
\begin{defn}\ollabel{defnomega}
لتكن~$I$ أي مجموعة تمنحنا إياها مسلّمة اللانهاية. ولتكن~$s$ الدالة
$s(x) = x \cup \{x\}$. ولتكن
$\omega = \closureofunder{s}{\emptyset}$. نسمي عناصر~$\omega$
\emph{الأعداد الطبيعية}، ونقول إن~$n$ ناتج تطبيق~$s$ على~$\emptyset$
$n$ مرة.
\end{defn}

يمكنك الآن الرجوع إلى الوراء والتحقق من أن المجموعة الموسومة «$n$»
قبل بضع فقرات ستُعامل \emph{بوصفها} العدد~$n$.

سنناقش أهمية هذا الاصطلاح في~\olref[sth][z][nat]{sec}. أما الآن فهو
يمكّننا من إثبات نتيجة حدسية:

\begin{prop}\ollabel{naturalnumbersarentinfinite}
لا عدد طبيعي لا نهائي على معنى ددكند.
\end{prop}

\begin{proof}
البرهان بالاستقراء، أي
\olref[sfr][infinite][induction]{thm:dedinfiniteinduction}. ومن الواضح
أن~$0 = \emptyset$ ليس لا نهائيًا على معنى ددكند. وفي خطوة الاستقراء
سنثبت المقابل بالنقيض: إذا كان~$s(n)$، على نحو محال، لا نهائيًا على
معنى ددكند، كان~$n$ لا نهائيًا على معنى ددكند.

فافترض أن~$s(n)$ لا نهائي على معنى ددكند، أي يوجد !!{injection}~$f$
بحيث~$\ran{f}\subsetneq \dom{f} = s(n) = n \cup \{n\}$. وهناك حالتان
ينبغي النظر فيهما.

\emph{الحالة 1: $n \notin \ran{f}$.} إذن~$\ran{f} \subseteq n$،
و$f(n) \in n$. ولتكن~$g = \funrestrictionto{f}{n}$؛ عندئذ
$\ran{g} = \ran{f} \setminus \{f(n)\} \subsetneq n = \dom{g}$.
ومن ثم يكون~$n$ لا نهائيًا على معنى ددكند.

\emph{الحالة 2: $n \in \ran{f}$.} ثبّت
$m \in \dom{f} \setminus \ran{f}$، وعرّف دالة~$h$ مجالها
$s(n) = n \cup \{n\}$:
\[
h(x) = 
	\begin{cases}
		f(x) & \text{إذا كان }f(x) \neq n\\
		m & \text{إذا كان }f(x)=n
	\end{cases}
\]
إذن تتفق~$h$ و$f$ في كل موضع، إلا أن
$h(f^{-1}(n)) = m \neq n = f(f^{-1}(n))$. ولما كانت~$f$
!!a{injection}، كان~$n \notin \ran{h}$؛ وكذلك
$\ran{h} \subsetneq \dom{h} = s(n)$. ويكون~$n$ الآن لا نهائيًا
على معنى ددكند باستعمال حجة الحالة~1.
\end{proof}

ومع ذلك يبقى السؤال عن كيفية \emph{تبرير} مسلّمة اللانهاية. والجواب
القصير هو أننا سنحتاج إلى إضافة مبدأ آخر إلى الرواية التي كنا نقدمها.
وهذا المبدأ هو:
\begin{enumerate}
	\item[] \stagesinf. توجد مرحلة لا نهائية. أي توجد مرحلة (أ) ليست المرحلة الأولى، و(ب) تسبقها بعض المراحل، لكنها (ج) ليس لها سابق مباشر.
\end{enumerate}
وتنتج مسلّمة اللانهاية مباشرة من هذا المبدأ. فنحن نعلم أن العدد
الطبيعي~$n$ يتكون في المرحلة~$n$. ولذلك تتكون المجموعة~$\omega$ في
أول مرحلة لا نهائية، وتكون~$\omega$ نفسها شاهدًا على مسلّمة اللانهاية.

غير أن هذا يعيدنا ببساطة إلى سؤال كيفية تبرير~\stagesinf. وكما في
\stagessucc، لم يكن هذا جزءًا صريحًا من الرواية التي قدمناها عن الهرم
التراكمي-التكراري. وأكثر من ذلك، لا يجبرنا شيء في فكرة الهرم التكراري
نفسها، حيث تتكون المجموعات مرحلة بعد مرحلة، على الاعتقاد بأن العملية
تتضمن مرحلة \emph{لا نهائية}. ويبدو منسجمًا تمامًا أن نتصور المراحل
مرتبة كالأعداد الطبيعية.

لكن هذا ينشئ مشكلة واضحة. فقد نظرنا في
\olref[sfr][infinite][dedekindsproof]{sec} في «برهان» ددكند على وجود
مجموعة (من الأفكار) لا نهائية على معنى ددكند. وربما لم تجده مقنعًا
جدًا. لكن إذا لم يكن~\stagesinf{} «مفروضًا علينا» بالتصور التكراري
للمجموعة (أو بـ«قوانين الفكر»)، فإننا نظل بلا تبرير داخلي للادعاء بوجود
مجموعة لا نهائية على معنى ددكند.

وثمة بالطبع ما يمكن قوله أكثر بكثير هنا. لكن يُرجى أنك بلغت الآن موضعًا
يمكنك فيه بدء التفكير فيما قد \emph{يلزم} لتبرير مسلّمة (أو مبدأ).
وفيما يأتي سنأخذ~\stagesinf{} مسلّمًا به ببساطة.

\end{document}
